\documentclass[10pt,a4paper]{amsart}
\usepackage[margin=19mm]{geometry}
\usepackage{amsmath,amssymb,mathtools}
\usepackage[scr=rsfs,cal=euler]{mathalfa}
\usepackage{xfrac}
\usepackage[unicode,hidelinks,bookmarks=false]{hyperref}
\newcommand{\ie}{i.e.\ }
\newcommand{\cf}{cf.\ }
\newcommand{\resp}{respectively\ }
\newcommand{\Subsection}{\subsection}
\providecommand{\nobreakdash}{\nobreak\hbox{-}}
\setlength{\emergencystretch}{2em}
\multlinegap0pt
\usepackage{amssymb}
\usepackage{enumerate}
\usepackage[shortlabels]{enumitem}
\newtheorem{thm}{Theorem}
\newtheorem{prop}{Proposition}
\newcommand{\Ric}{\mathrm{Ric}}
\newcommand{\Vol}{\mathrm{Vol}}
\newcommand{\bC}{\mathbb{C}}
\newcommand{\cK}{\mathcal{K}}
\newcommand{\idd}{i\partial\bar\partial}
\newcommand{\loc}{\mathrm{loc}}
\newcommand{\supp}{\mathrm{supp}}
\newcommand{\trdeg}{\mathrm{trdeg}}
\title{Revised Demailly's Affineness Criterion and Algebraization of Entire Grauert Tubes}
\author{Kyobeom Song}
\date{Source: arXiv:2605.05582v2 (CC BY 4.0)}
\begin{document}
\maketitle

\noindent\emph{Source and licence.} Revised Demailly's Affineness Criterion and Algebraization of Entire Grauert Tubes. Version arXiv:2605.05582v2 (CC BY 4.0), distributed by arXiv under the Creative Commons Attribution 4.0 International licence, http://creativecommons.org/licenses/by/4.0/, as recorded in the arXiv licence metadata for this exact version and shown on the abstract page https://arxiv.org/abs/2605.05582v2. Full licence notice: \texttt{licenses/arxiv-2605.05582-CC-BY-4.0.txt}.\par\smallskip

\noindent\textit{Editorial scope.} Counted unit: the article's prop jet (source0.tex lines 404--422). The article's complete original proof of it is delivered with it (source0.tex lines 424--531, one \texttt{proof} environment of 108 lines). No mathematics is written, repaired or re-derived: the statement and the proof are the article's own lines, in the article's own notation, and no retained line is substituted or re-flowed.  Retained uncounted as the source-local context the proof needs: lines 322--332.  Nothing was omitted from any retained span, so the package carries no omission record.  No retained line contains a citation, so the wrapper carries no bibliography and no bibliography entry.  No retained line contains a figure environment, an image-inclusion command or drawing code, so no image is delivered and the package declares no asset dependency.  Editorial formatting only, in addition: an AMS \texttt{amsart} preamble that restates the article's own declarations and macros used by the retained lines, namely the environments \texttt{thm}, \texttt{prop} on separate counters, together with the standard packages those lines need.  The retained environments print in this document as Theorem 1, Proposition 1 in order of appearance, whereas the article prints the same items with the numbering of its own full document.  The complete native source of the article as distributed in the arXiv e-print for this exact version is bundled unchanged as \texttt{sources/199967/source0.tex}.
\begin{thm}\label{trdeg}
Let $X$ be an $n$-dimensional Stein manifold with finite Monge--Amp\`ere volume
\[
\Vol(X):=\int_X (dd^c\phi)^n < \infty.
\]
Then the following hold:
\begin{enumerate}
\item $0\le \trdeg_{\bC}\cK_\phi(X)\le n$.
\item If $\trdeg_{\bC}\cK_\phi(X)=n$, then $\cK_\phi(X)$ is a finite type extension of $\bC$.
\end{enumerate}
\end{thm}

\begin{prop}\label{jet}
The following statements hold.
\begin{enumerate}
\item[(a)] Given finitely many points $p_1,\dots,p_k\notin S':=\{x~|~e^{-\psi}\notin L^1_\loc(x)\}\subset S$, and prescribed $s_i$-jets at each
$p_i$, there exists $f\in A_\phi^0(X)$ realizing the prescribed $s_i$-jet at each $p_i$.

\item[(b)] Let $f_1,\dots,f_n\in A_\phi^0(X)$ satisfy
\[
\bigwedge_{i=1}^n df_i(p)\neq 0
\qquad\text{for some }p\notin S.
\]
Then there exists $b_p\in A_\phi^0(X)$ such that $b_p(p)=1$ and
\[
\{\textstyle\bigwedge_{i=1}^n df_i=0\}\subset \{b_p=0\}.
\]

\item[(c)] $\trdeg_\mathbb{C}\cK_\phi(X)=n$, and $\cK_\phi(X)$ is a finite type extension of $\mathbb{C}$.
\end{enumerate}
\end{prop}

\begin{proof}
\noindent (a)
First, note that $S$ and $S'$ are closed. Hence, we can fix finitely many points $p_1,\dots,p_k$ and choose local coordinate systems
$z^{(1)},\dots,z^{(k)}$ near $p_1,\dots,p_k$, respectively, such that these coordinate neighborhoods are disjoint
from $S'$.
Let $\chi_i$ be a cutoff function supported in the coordinate neighborhood of $z^{(i)}$ and identically $1$
near $p_i$.
For each $p_i$, choose a polynomial $P_i$ of degree $s_i$ in $z^{(i)}$ whose $s_i$-jet at $p_i$ is the desired one.
Set
\[
h \;:=\; \sum_i \chi_i P_i,\qquad u\;:=\;\bar\partial h,\qquad s\;:=\;\max_i s_i.
\]
Choose $C>0$ sufficiently large so that
\[
\eta \;:=\; C\phi + (n+s)\sum_i \chi_i \log|z^{(i)}|^2
\]
is plurisubharmonic. Define
\[
\tau \;:=\; \psi+\eta+2\log(1+e^\phi).
\]
Then, since $\Ric(\beta)+\idd\psi\ge 0$ and $dd^c\log(1+e^\phi)\ge\frac{1}{(1+e^\phi)^2}\beta$, one obtains
\[
\idd\tau+\Ric(\beta)\;\ge\;\frac{1}{(1+e^\phi)^2}\,\beta.
\]
Hence we may take $\lambda=(1+e^\phi)^{-2}$.
Since $u:=\bar \partial h$ vanishes near $p_1,\cdots,p_k$ and $\supp(u)$ is relatively compact, while $e^{-\psi}$ is locally integrable on $\supp(u)$, we have
\[
\int_X \lambda^{-1}|u|_\beta^2 e^{-\tau}\beta^n <\infty.
\]
By the lemma there exists $g$ with $\bar\partial g=u$ and $\int_X |g|^2 e^{-\tau}\beta^n<\infty$.
Moreover, since $\psi\le A\phi+B$, $\eta\le C\phi+C'$ for some $C'>0$, and
$2\log(1+e^\phi)\le 2\phi+2\log 2$, there exist constants $A',B'$ such that
\[
e^{-\tau}\;\ge\; e^{-(A'\phi+B')}.
\]
Thus
\[
\int_X |g|^2 e^{-(A'\phi+B')}\beta^n <\infty,
\]
so $g\in L_\phi^0(X)$.
Finally, since $\bar\partial h=u$ and $\bar\partial g=u$, we have $\bar\partial(h-g)=0$, hence
\[
f\;:=\;h-g \in \mathcal{O}(X),
\qquad\text{and in particular } f\in A_\phi^0(X).
\]
The finiteness
\[
\int_X |g|^2 e^{-\tau}\beta^n
=\int_X |g|^2 e^{-\psi-C\phi}(1+e^\phi)^{-2}
\Bigl(\prod_i |z^{(i)}|^{-2(n+s)\chi_i}\Bigr)\,\beta^n \;<\;\infty
\]
forces $g$ to have no nonzero jets of order $\le s$ at each $p_i$.
Therefore $f=h-g$ has, at each $p_i$, the prescribed jet (matching that of $h$ up to order $s_i$).

\medskip
\noindent (b)
We apply the same argument as (a) with a slight modification.
Choose a local coordinate system $z$ near $p$ such that the neighborhood is disjoint from
$S$ and from $\{\bigwedge_{i=1}^n df_i=0\}$.
Let $\chi$ be a cutoff supported in this coordinate neighborhood and identically $1$ near $p$.
Set $h:=\chi$ and
\[
\eta \;:=\; C\phi + n\,\chi \log|z|^2,
\]
for some $C>0$ chosen so that $\eta$ is psh.
Define
\[
\tau \;:=\; 2\psi+\eta+\log\bigl|\textstyle\bigwedge_{i=1}^n df_i\bigr|_\beta^2 + 2\log(1+e^\phi).
\]
Let $[Z]$ denote the zero divisor of $\bigwedge_{i=1}^n df_i$ (viewed as a $(1,1)$-current).
Then
\[
\idd\tau+\Ric(\beta)
=2\idd\psi+2\Ric(\beta)+2\pi[Z]+\idd\eta+dd^c\log(1+e^\phi)
\;\ge\;\frac{1}{(1+e^\phi)^2}\,\beta.
\]
As in the proof of (a), $\bar\partial h$ vanishes near $p$ and $\supp(\bar\partial h)$ is relatively compact. Moreover, on $\supp(\bar\partial h)$, $e^{-2\psi}$ is locally integrable and $\log\bigl|\textstyle\bigwedge_{i=1}^n df_i\bigr|_\beta^2$ is finite.  Therefore, taking again $\lambda=(1+e^\phi)^{-2}$ yields
\[
\int_X \lambda^{-1}|\bar\partial h|_\beta^2 e^{-\tau}\beta^n<\infty.
\]
By the lemma there exists $g$ with $\bar\partial g=\bar\partial h$ and
\[
\int_X |g|^2 e^{-\tau}\beta^n
=\int_X |g|^2 e^{-2\psi-C\phi}\,|z|^{-2n\chi}\,
\bigl|\textstyle\bigwedge_{i=1}^n df_i\bigr|_\beta^{-2}\,(1+e^\phi)^{-2}\,\beta^n
\;<\;\infty.
\]
The finiteness of this integral, despite the factors
$\bigl|\bigwedge_{i=1}^n df_i\bigr|_\beta^{-2}$ and $|z|^{-2n\chi}$, forces $g$ to vanish on
$\{\bigwedge_{i=1}^n df_i=0\}$ and at $p$.
Define $b_p:=h-g$. Then $\bar\partial b_p=0$, so $b_p$ is holomorphic, $b_p(p)=1$, and
$\{\bigwedge_{i=1}^n df_i=0\}\subset\{b_p=0\}$.
Moreover, since $\psi\le A\phi+B$, $\eta\le C\phi+C'$ for some $C'>0$, and
$2\log(1+e^\phi)\le 2\phi+2\log 2$, there exist $A',B'$ such that
\[
e^{-\tau}\;\ge\; e^{-(A'\phi+B')}\,\bigl|\textstyle\bigwedge_{i=1}^n df_i\bigr|_\beta^{-2}.
\]
Hence $|g|\cdot\bigl|\bigwedge_{i=1}^n df_i\bigr|_\beta^{-1}\in L_\phi^2(X)$, and therefore
\[
|g|\;\le\;\Bigl(|g|\cdot\bigl|\textstyle\bigwedge_{i=1}^n df_i\bigr|_\beta^{-1}\Bigr)
\prod_{i=1}^n |df_i|_\beta \in L_\phi^0(X),
\]
which implies $b_p=h-g\in A_\phi^0(X)$.

\medskip
\noindent (c)
By (a), we can construct $n$ functions $f_1, \cdots, f_n$ which separate all $1$-jets at a point $p\notin S$. These functions must be algebraically independent, hence $\trdeg_{\mathbb{C}} \cK_\phi(X) \ge n$. By Theorem~\ref{trdeg}, it follows that $\trdeg_{\mathbb{C}} \cK_\phi(X) = n$, and $\cK_\phi(X)$ is a finite type extension of $\mathbb{C}$.
\end{proof}

\end{document}
